\documentclass[UTF8,a4paper,11pt]{ctexart}

\usepackage[left=25mm,right=25mm,top=23mm,bottom=25mm]{geometry}
\usepackage{amsmath,amssymb,amsthm}
\usepackage{xcolor}
\usepackage{booktabs,tabularx}
\usepackage{enumitem}
\usepackage{fancyhdr}
\usepackage[hidelinks]{hyperref}

\definecolor{ecnu}{HTML}{1F5A94}
\definecolor{softblue}{HTML}{EEF5FB}
\setlength{\parindent}{2em}
\setlength{\parskip}{2pt}
\setlist{itemsep=2pt,topsep=3pt}
\pagestyle{fancy}
\fancyhf{}
\lhead{\color{ecnu}\textbf{数据结构}}
\rhead{作业二 · Your Name}
\cfoot{\thepage}
\renewcommand{\headrulewidth}{0.5pt}
\renewcommand{\headrule}{\hbox to\headwidth{\color{ecnu}\leaders\hrule height \headrulewidth\hfill}}

\newcounter{problem}
\newenvironment{problem}[1][]
  {\refstepcounter{problem}\par\vspace{8pt}\noindent\colorbox{ecnu}{\parbox{\dimexpr\linewidth-2\fboxsep}{\color{white}\textbf{Problem \theproblem}\if\relax\detokenize{#1}\relax\else\quad #1\fi}}\par\smallskip\ignorespaces}
  {\par}
\newenvironment{solution}
  {\par\noindent\color{ecnu}\textbf{Solution}\quad\color{black}\ignorespaces}
  {\par}
\newtheorem{theorem}{定理}

\begin{document}

\begin{center}
  {\zihao{1}\color{ecnu}\bfseries 数据结构与算法作业二\par}
  \vspace{5pt}
  {\large Graphs, recursion, and invariants\par}
  \vspace{5pt}
  {课程：数据结构\quad 姓名：Your Name\quad 学号：00000000\par}
\end{center}

\noindent\colorbox{softblue}{\parbox{\dimexpr\linewidth-2\fboxsep}{\textbf{要求：} 写出关键步骤和复杂度；可以使用英文或中文解释，但请保持符号定义一致。}}

\begin{problem}[递归与复杂度]
设 $T(n)=2T(n/2)+n$，其中 $T(1)=1$。使用递归树说明 $T(n)$ 的渐近复杂度，并解释每一层的工作量。
\begin{solution}
第 $i$ 层有 $2^i$ 个子问题，每个子问题的非递归工作为 $n/2^i$，因此这一层的总工作量为 $n$。树高为 $\log_2 n$，叶节点总工作量为 $n$，所以
\[
  T(n)=\Theta(n\log n).
\]
这也与主定理中 $a=2,b=2,f(n)=n$ 的结论一致。
\end{solution}
\end{problem}

\begin{problem}[最短路径]
给出一个使用优先队列的 Dijkstra 算法框架，并说明非负边权是算法正确性的必要条件。
\begin{solution}
将源点的距离设为零，其余距离设为无穷大。每次从优先队列取出当前距离最小的未确定顶点 $u$，再对每条边 $(u,v)$ 执行松弛：若 $d(v)>d(u)+w(u,v)$，则更新 $d(v)$ 并重新加入队列。因为边权非负，已经取出的最小标签不会被未来路径改善；若存在负边，这个单调性可能失效。
\end{solution}
\end{problem}

\begin{problem}[证明不变量]
为下面的循环写出一个不变量，并用三句话证明它保持成立：
\[
\texttt{for i = 1,\ldots,n:}\qquad s \leftarrow s+a_i.
\]
\end{problem}

\begin{theorem}
若循环开始前 $s=0$，循环结束后 $s=\sum_{i=1}^{n}a_i$。
\end{theorem}

\end{document}
